\chapter{Судалгааны сэдвийн онол, өнөөгийн түвшин}
\section{Онолын хэсэг}

\subsection{Био-метрик танилт (Biometric Recognition)}
Био-метрик (biometric) танилт нь хүний нүүр, хурууны хээ, дуу хоолой, нүдний торлог гэх мэт өөрчлөхөд бэрх шинжүүдийг ашиглан таних эсвэл баталгаажуулах асуудлыг шийддэг. 
Танилтын чанарыг нарийвчлал (accuracy), үнэн позитивын хувь (TPR), 
худал позитивын хувь (FPR) болон ROC/EER зэрэг хэмжигдэхүүнээр үнэлнэ.

\subsection{Нүүр илрүүлэлт (Face Detection)}

Нүүр илрүүлэлтийн эхний шат нь дүрс эсвэл видео кадрт орсон хүний нүүрийг илрүүлэх (face detection) үйл явц юм. Энэ шатанд Haar Cascade, HOG+SVM, MTCNN, RetinaFace зэрэг алгоритмуудыг ашиглах боломжтой. 

RetinaFace нь нэг шатлалт илрүүлэгч (one-stage detector) бөгөөд нэг forward pass үйлдлээр олон даалгаврын (multi-task) гаралт үүсгэдэг. Үүнд нүүрний байрлалын хүрээ (bounding box), нүүрний landmark цэгүүд болон итгэлцлийн үнэлгээ (confidence score) багтана.

Энэхүү архитектур нь жижиг хэмжээтэй нүүр, өнцгийн өөрчлөлттэй нөхцөлд илрүүлэлтийн тогтвортой байдлыг сайжруулдаг. Мөн landmark цэгүүдийг ашигласнаар нүүрийг тэгшлэх (alignment) болон тайрах (crop) шатны чанарыг нэмэгдүүлдэг. Иймээс бодит орчны ирц бүртгэлийн системд RetinaFace нь практик хэрэглээнд тохиромжтой сонголт гэж үзэж болно \cite{deng2019retinaface}.

\subsection{Нүүрийн хэв шинж ба эмбеддинг (Feature Extraction and Embedding)}
Нүүрийг пикселийн түвшний дүрсийн мэдээллээс ангилалд тохиромжтой тоон төлөөлөл болгон хувиргахын тулд эмбеддинг буюу нүүрний ялгаатай онцлогийг агуулсан вектор шаардлагатай байдаг. Энэхүү баримт бичигт нүүрний эмбеддинг-г $f(x) \in \mathbb{R}^{128}$ хэмжээст вектор гэж авч үзнэ. FaceNet загвар нь уг векторыг triplet loss функц ашиглан сургадаг бөгөөд төстэй царайтай хүмүүсийн эмбеддингүүдийг векторын орон зайд ойр байрлуулж, төсгүй царайтай хүмүүсийн эмбеддингүүдийг хооронд нь хол байрлуулах байдлаар ялгаатай төлөөлөл үүсгэдэг \cite{schroff2015facenet}.

Triplet loss функцийн математик илэрхийлэл нь дараах хэлбэртэй байна:
\[
L = \max \left(0, \lVert f(a)-f(p)\rVert_2^2 - \lVert f(a)-f(n)\rVert_2^2 + \alpha \right)
\]

Энд $a$ нь anchor жишээ, $p$ нь anchor-той ижил ангилалд хамаарах positive жишээ, $n$ нь өөр ангилалд хамаарах negative жишээ, $\alpha$ нь ангиллуудын хоорондын зайн доод хязгаарыг тогтоох margin параметр болно.

Илрүүлэлтийн шатны чанар буурах нь эмбеддингийн чанарт шууд нөлөөлдөг. Тухайлбал, буруу илрүүлэлт эсвэл тэгшитгэл нь чанар муутай тайралт үүсгэж, улмаар шуугиант эмбеддинг бий болгоно. Үүний үр дүнд ангилалтын алдааны түвшин өсөх эрсдэлтэй.

Иймээс илрүүлэлт болон эмбеддинг шатны чанарын уялдаа нь системийн нийт нарийвчлалын гол тодорхойлогч хүчин зүйл болдог.

\subsection{Эмбеддингийн тогтвортой байдалд нөлөөлөх хүчин зүйлс (Robustness Factors)}
Бодит орчинд ижил хүний нүүрний дүрсүүд нь гэрэлтүүлэг, өнцөг, хөдөлгөөний бүдгэрэл, халхлалт зэрэг шалтгаанаар ялгаатай тархалттай болдог. Үүний улмаас ижил хүний эмбеддингүүдийн зай өсөж, өөр хүмүүсийн эмбеддингүүдтэй огтлолцох эрсдэл нэмэгддэг.

Энэ судалгааны хүрээнд дараах хүчин зүйлс онцгой нөлөөтэй гэж үзсэн.
\begin{itemize}
\item Гэрэлтүүлэг: Нүүрний локал контраст багасах эсвэл хэт ихсэхэд ялгах шинж тэмдгүүдийн тогтвортой байдал буурна.
\item Өнцөг (pose): Эгц урдаас бус өнцгөөс геометрийн хувирал нэмэгдэж, алдаа өснө.
\item Хөдөлгөөний бүдгэрэл: Нүүрний нарийн ширхэгтэй бүтэц бүдгэрч, ангилалд чухал ялгах мэдээлэл буурна.
\item Хэсэгчилсэн халхлалт: Нүд, хамар, ам орчмын түлхүүр бүсүүд хаагдахад эмбеддинг векторын зөв төлөөлөх чадвар буурна.
\end{itemize}

Дээрх онолын нөлөөллүүд нь судалгааны үр дүн бүлгийн Зураг~\ref{fig:error-cases-annotated}-д ажиглагдсан алдаатай нийцэж байна. Тухайлбал, ижил төстэй өнцөгтэй хоёр тохиолдолд нарны шууд тусгал нэмэгдэх үед танилтын итгэлцүүр буурч \texttt{танигдаагүй} төлөв рүү шилжсэн нь гэрэлтүүлгийн өөрчлөлт эмбеддингийн тогтвортой байдалд хүчтэй нөлөөлж байгааг харуулсан. Мөн FP тохиолдол ажиглагдсан нь зарим нөхцөлд ангиллын зааг ойртож болохыг баталж байна.

\subsection{Ангилагч(Classifier)}
Эмбеддинг үүссэний дараа тухайн хүнийг танихын тулд ангилагч ашиглана. Ангилагч нь эмбеддинг векторуудыг авч, тухайн жишээ аль ангилалд хамаарахыг тодорхойлно.

Танилтын шатанд cosine similarity хэмжүүрийг шууд ашиглах боломжтой:
\[
\cos(\theta)=\frac{u\cdot v}{\lVert u\rVert\,\lVert v\rVert}
\]
энд $u,v \in \mathbb{R}^{128}$ байна. Cosine утга 1-д ойртох тусам хоёр эмбеддинг векторын чиглэл илүү төстэй, өөрөөр хэлбэл нэг ижил хүний төлөөлөл байх магадлал өндөр гэж үзнэ.

Мөн k-Nearest Neighbors (KNN) алгоритмыг ашиглаж болно. KNN нь шинэ эмбеддингэд хамгийн ойр орших $k$ хөршийг cosine эсвэл Euclidean зайгаар тодорхойлж, олонхийн саналын зарчмаар ангиллыг сонгодог. Энэ арга нь параметрчилсэн сургалт шаарддаггүй (non-parametric) тул шинэ хүн нэмэх үед загварыг дахин сургах шаардлагагүй, зөвхөн шинэ эмбеддингийг санах ойд нэмэхэд хангалттай. 
Гэвч өгөгдлийн хэмжээ ихсэхэд хайлтын хугацаа нэмэгдэх сул талтай бөгөөд үүнийг индексжүүлэлтийн арга (жишээлбэл, хамгийн ойрын хөрш хайлт) ашиглан бууруулах боломжтой \cite{chen2019face}.

\subsection{Илрүүлэлтийн аргуудын харьцуулалт}
Нүүр илрүүлэлтийн шатанд ашиглагдах алгоритмууд нь өгөгдлийн чанар, гэрэлтүүлгийн нөхцөл, камерын өнцөг, мөн бодит цагийн (real-time) гүйцэтгэлийн шаардлагаас хамааран өөр өөр давуу болон сул талтай байдаг. Иймээс практик хэрэглээнд тохирсон алгоритмыг сонгохын тулд тэдгээрийг харьцуулан үнэлэх нь зүйтэй.

\begin{table}[H]
\centering
\caption{Нүүр илрүүлэлтийн аргуудын харьцуулалт}
\begin{tabularx}{\textwidth}{|l|X|X|}
\hline
\textbf{Арга} & \textbf{Давуу тал} & \textbf{Сул тал} \\
\hline
Haar Cascade & Хурдан, хөнгөн, CPU дээр хэрэгжүүлэхэд хялбар & Өнцөг болон гэрэлтүүлгийн өөрчлөлтөнд тогтвортой байдал буурдаг \\
\hline
HOG+SVM & Түгээмэл хэрэглэгддэг арга, хэрэгжүүлэхэд энгийн & Жижиг хэмжээтэй болон төвөгтэй нөхцөлтэй нүүрэн дээр нарийвчлал буурдаг \\
\hline
MTCNN & Bounding box болон landmark цэгүүдийг хамтад нь илрүүлдэг, олон нөхцөлд нарийвчлал сайн & Тооцооллын зардал Haar болон HOG аргуудаас өндөр \\
\hline
RetinaFace & Хүнд нөхцөлд илрүүлэлтийн чанар өндөр, жижиг нүүрэнд үр дүн сайтай & Илүү их тооцооллын болон санах ойн нөөц шаарддаг \\
\hline
\end{tabularx}
\end{table}

\subsection{Эмбеддинг аргуудын харьцуулалт}

Нүүрний эмбеддинг загваруудын ялгаа нь ашиглагдах алдагдлын функц (loss function), эмбеддинг векторуудын хоорондын ялгаралтыг нэмэгдүүлэх механизм, мөн бодит орчны нөхцөлд үзүүлэх тогтвортой байдлаар тодорхойлогдоно.

\begin{table}[H]
\centering
\caption{Нүүрний эмбеддинг аргуудын харьцуулалт}
\begin{tabularx}{\textwidth}{|l|X|X|}
\hline
\textbf{Арга} & \textbf{Давуу тал} & \textbf{Сул тал} \\
\hline
FaceNet (Triplet loss) & Ерөнхий зориулалтын эмбеддинг үүсгэдэг, open-set танилтад (урьдчилан харагдаагүй ангилалтай ажиллах) өргөн хэрэглэгддэг & Triplet mining стратеги болон сургалтын өгөгдлийн чанарт мэдрэмтгий \\
\hline
SphereFace (Angular margin) & Анги хоорондын өнцгийн зайг нэмэгдүүлж ялгах чадварыг сайжруулдаг & Сургалтын явцад convergence тогтвортой бус байх магадлалтай \\
\hline
ArcFace (Additive angular margin) & Ялгаралт өндөртэй эмбеддинг үүсгэж, танилтын нарийвчлал сайжруулдаг & Их хэмжээний өгөгдөл болон тооцооллын нөөц шаарддаг \\
\hline
CosFace (Cosine margin) & Cosine зайд суурилсан margin ашиглан анги хоорондын ялгаралтыг нэмэгдүүлдэг & Гипер-параметрийн тохиргоо буруу бол гүйцэтгэл буурах эрсдэлтэй \\
\hline
\end{tabularx}
\end{table}

\subsection{Ангиллын аргууд ба сургалтын онцлог}
Ижил эмбеддинг орон зайд ангиллын шатанд ерөнхийдөө хоёр үндсэн бүлгийн арга хэрэглэгддэг. 

Нэгдүгээрт, параметрчилсэн сургалт шаардахгүй (non-parametric) аргууд болох KNN болон cosine-threshold арга нь шинэ хүн нэмэх үед эмбеддинг санг шинэчлэхэд хангалттай бөгөөд загварыг дахин сургах шаардлагагүй.

Хоёрдугаарт, параметрчилсэн сургалт шаарддаг (parametric) аргууд болох SVM болон Softmax/ArcFace ангилагч нь тодорхой сургалтын өгөгдлөөр параметрүүдийг сургаж, ангийн тоо нэмэгдэх үед загварыг дахин сургах шаардлага үүсгэж болзошгүй байдаг.

\begin{table}[H]
\centering
\caption{Ангиллын аргуудын харьцуулалт ба сургалтын онцлог}
\begin{tabularx}{\textwidth}{|l|X|X|X|}
\hline
\textbf{Арга} & \textbf{Сургалтын арга} & \textbf{Давуу тал} & \textbf{Сул тал} \\
\hline
Cosine similarity & Сургалтгүй (босго параметр тохируулна) & Хурдан, хэрэгжүүлэхэд энгийн, open-set нөхцөлд уян хатан & Босго буруу тохирвол худал эерэг (false positive) болон худал сөрөг (false negative) алдаа нэмэгдэнэ \\
\hline
KNN & Параметрчилсэн сургалтгүй, эмбеддинг санг шинэчилнэ & Шинэ хэрэглэгч нэмэхэд дахин сургах шаардлагагүй & Сангийн хэмжээ ихсэхэд хайлтын хугацаа өснө \\
\hline
SVM & Хяналттай сургалт (supervised learning) & Ангиллын хилийг оновчтой тогтоож, дундаж хэмжээний өгөгдөлд сайн гүйцэтгэл үзүүлдэг & Анги нэмэгдэхэд дахин сургах шаардлагатай \\
\hline
Softmax/ArcFace head & Deep classifier сургалт (cross-entropy + margin loss) & Closed-set нөхцөлд өндөр нарийвчлал гаргах боломжтой & Тооцооллын нөөц их шаарддаг, шинэ анги нэмэхэд уян хатан бус \\
\hline
\end{tabularx}
\end{table}

\subsection{Үнэлгээний хэмжүүрүүдийн онолын үндэс}

Царай танилтын системийн гүйцэтгэлийг бодитой үнэлэхийн тулд confusion matrix-д суурилсан ангиллын хэмжүүрүүдийг ашиглах нь зүйтэй. Энд $TP$ (True Positive), $FP$ (False Positive), $TN$ (True Negative), $FN$ (False Negative) утгууд нь ангиллын үр дүнгийн бүтцийг тодорхойлж, системийн алдааны шинж чанарыг тайлбарлана.

\begin{table}[H]
\centering
\caption{Үндсэн ангиллын хэмжүүрүүдийн тайлбар}
\label{tab:metric-glossary-theory}
\begin{tabularx}{\textwidth}{|l|X|}
\hline
\textbf{Хэмжүүр} & \textbf{Тайлбар} \\
\hline
TP & Бодитоор тухайн оюутан байсан бөгөөд систем зөв таньсан тохиолдол. \\
\hline
FP & Бодитоор тухайн оюутан биш боловч систем андуурч таньсан (буруу эерэг шийдвэр) тохиолдол. \\
\hline
TN & Бодитоор тухайн оюутан байгаагүй бөгөөд систем зөвөөр тухайн ангилалд оруулаагүй тохиолдол. \\
\hline
FN & Бодитоор тухайн оюутан байсан боловч систем таньж чадаагүй эсвэл \texttt{unknown} ангилалд шилжүүлсэн тохиолдол. \\
\hline
\end{tabularx}
\end{table}

Нийт зөв ангилсан хувийг Accuracy хэмжүүрээр илэрхийлнэ:

\[
\text{Accuracy} = \frac{TP + TN}{TP + TN + FP + FN} \times 100
\]

Accuracy нь системийн ерөнхий гүйцэтгэлийг харуулах боловч анги хоорондын тэнцвэргүй байдалтай өгөгдөлд хангалттай мэдээлэл өгөхгүй байж болно.

Системийн эерэг шийдвэрийн найдвартай байдлыг Precision-ээр үнэлнэ:

\[
\text{Precision} = \frac{TP}{TP + FP} \times 100
\]

Precision өндөр байх нь буруу эерэг (FP) шийдвэр бага байгааг илтгэнэ.

TP тохиолдлыг зөв таних чадварыг Recall-ээр үнэлнэ:

\[
\text{Recall} = \frac{TP}{TP + FN} \times 100
\]

Recall өндөр байх нь бодит эерэг тохиолдлуудыг алдалгүй таньж байгааг илэрхийлнэ.

Precision болон Recall-ийн тэнцвэржүүлсэн нийлмэл үзүүлэлт нь F1-score болно:

\[
\text{F1-score} = \frac{2 \times \text{Precision} \times \text{Recall}}{\text{Precision} + \text{Recall}}
\]

F1-score нь ангиллын алдааны хоёр төрлийг (FP болон FN) зэрэг харгалзан үздэг тул тэнцвэртэй үнэлгээ болдог.

Мөн open-set нөхцөлд \texttt{unknown} ангилал нь $FP$ шийдвэрийг бууруулах чухал механизм болдог. Гэвч босго параметр хэт хатуу тохируулагдсан тохиолдолд $FN$ нэмэгдэх эрсдэлтэй тул $\tau$ төрлийн босгуудыг өгөгдөлд тулгуурлан тэнцвэржүүлэх шаардлагатай.
\section{Ижил төрлийн судалгаа}

\subsection{Гадаадын судалгаа}

Олон улсын хэмжээнд камерын дүрсэнд суурилсан ирц бүртгэлийн системүүдийг өргөнөөр судалж, туршиж байна. Эдгээр системүүд нь ихэвчлэн царай таних технологи, машин сургалтын ангилагч болон бодит цагийн дүрс боловсруулалтын аргачлалд тулгуурладаг.

Жишээлбэл, FaceNet загварыг ашигласан судалгаанууд нь сурагчдын нүүрний эмбеддинг-г үүсгэж, cosine similarity хэмжүүрээр танилтыг гүйцэтгэх замаар өндөр нарийвчлалтай ирц бүртгэх боломжтойг харуулсан \cite{schroff2015facenet}. Мөн ArcFace болон SphereFace зэрэг margin-based softmax аргууд нь эмбеддинг орон зай дахь анги хоорондын ялгаралтыг нэмэгдүүлж, ангилалтын нарийвчлалыг сайжруулах боломжтойг тогтоосон \cite{deng2019arcface, liao2017sphereface}.

Ангиллын шатанд KNN, SVM зэрэг уламжлалт аргуудыг эмбеддинг дээр суурилан ашиглах судалгаа түгээмэл байна. KNN арга нь шинэ анги нэмэгдэх үед дахин сургалт шаардахгүй давуу талтай бол SVM нь ангиллын хилийг margin зарчмаар тодорхойлж, дундаж хэмжээний өгөгдөлд тогтвортой гүйцэтгэл үзүүлдэг \cite{chen2019face}.

Сүүлийн жилүүдэд MTCNN, RetinaFace зэрэг гүн сургалтын алгоритмуудыг ашиглан нүүр илрүүлэлтийн нарийвчлалыг сайжруулах судалгаа эрчимжсэн байна. Мөн FaceNet, ArcFace зэрэг эмбеддинг-д суурилсан загваруудыг ашиглан танилтын алдааг бууруулах чиглэл давамгайлж байна. Зарим судалгаанд KNN, SVM, cosine similarity зэрэг энгийн боловч үр дүнтэй ангилагчдыг ашиглан бодит орчинд хурдан ажиллах боломжийг судалсан байдаг \cite{luo2025opencv,devi2025svm}.

Үүнээс гадна edge computing болон IoT/embedded орчныг ашиглан ангид шууд тооцоолол хийх шийдлүүдийг дэвшүүлсэн нь сүлжээний хамаарлыг багасгаж, бодит хугацааны гүйцэтгэлийг сайжруулах боломжтойг тэмдэглэсэн \cite{yadav2021inbrowser,vu2024ble}.

Эдгээр судалгааны үр дүнгээс харахад царай танилтад суурилсан ирц бүртгэлийн систем нь зохистой алгоритмын сонголт, өгөгдлийн чанарын хангалттай нөхцөлд өндөр нарийвчлал, автоматжуулалтын түвшинг хангах боломжтойг харуулж байна.

\subsection{Дотоодын хэрэгжилт}

Монгол Улсад царай танилтад суурилсан ирц бүртгэлийн чиглэлээр диплом болон төслийн ажлууд сүүлийн жилүүдэд хийгдэж эхэлсэн нь судалгааны суурь бүрэлдэж буйг илтгэнэ. Жишээлбэл, ШУТИС-ийн цахим номын сангийн бүртгэлтэй ``Хиймэл оюун ухаанд суурилсан царай таньдаг ирц бүртгэлийн системийн хөгжүүлэлт'' (2024) болон ``Царай танилтаар бүртгэх систем'' (2025) зэрэг бакалаврын ажлууд нь нүүр илрүүлэлт, танилт болон ирцийн автомат бүртгэлийн прототип шийдлүүдийг дэвшүүлсэн байна \cite{must_enkhjin2024,must_uuganbayar2025}.

Гэвч эдгээр ажлууд нь ихэвчлэн туршилт, прототипийн түвшинд хязгаарлагдсан бөгөөд ерөнхий боловсролын болон их, дээд сургуулиудын хэмжээнд өргөн хэрэглээний тогтмол систем болон нэвтрээгүй, нэвтрүүлэлт эхлэл шатандаа байна. Иймд илүү өргөн хүрээнд туршилт хийх, алгоритмын нарийвчлалыг сайжруулах, хэрэглээний уян хатан байдлыг нэмэгдүүлэх чиглэлээр судалгаа, хөгжүүлэлт хийх шаардлагатай байна.
\subsection{Судалгааны онцлог}

Одоогийн зарим камерын суурьтай ирц бүртгэлийн системүүд нь бодит орчинд удаан хугацаанд тасралтгүй ажиллах үед гүйцэтгэлийн ачаалал нэмэгдэж, илрүүлэлт болон танилтын нарийвчлал буурах сул талтай. Мөн зарим систем нь зөвхөн анги танхимд нэвтрэх үеийн бүртгэлд төвлөрдөг бөгөөд хичээлийн явцад тасралтгүй хяналт хийх боломж хязгаарлагдмал байдаг.

Иймээс илрүүлэлт болон танилтын шатанд нарийвчлал сайтай алгоритмууд (жишээлбэл, MTCNN, FaceNet гэх мэт) ашиглаж, бодит орчны нөхцөлд тогтвортой, найдвартай ажиллах ирц бүртгэлийн системийн загварыг боловсруулах шаардлагатай байна. Дараагийн бүлэгт дээрх харьцуулалтын дүгнэлтэд тулгуурлан сонгосон архитектур, өгөгдөл боловсруулах урсгал болон үнэлгээний аргачлалыг дэлгэрэнгүй тайлбарлана.